\section{Convergence for the Laurent symbol algebras}\label{su:sec}

\providecommand{\fA}{\mathfrak A}
\providecommand{\hot}{\mathbin{\widehat\otimes}}
\newcommand{\Ssym}{\mathcal S}
\newcommand{\Gsym}{\mathcal G}
\newcommand{\Omh}{\Omega_{\mathrm h}}

This appendix proves (B10) for every \(n\ge1\): the order spectral
sequences of \(\Ads\) and \(\As\) converge, and the leading-symbol map
\(\Theta\) is an isomorphism (Theorem~\ref{su:thm:B10}). This closes the
gap described in Remarks~\ref{ha:rem:quoted} and~\ref{ha:rem:xadicfail}.

The argument is due to Lauter and Moroianu \cite{LauterMoroianu}, who
compute the Hochschild homology of the algebras of fibered cusp
operators. Their fibered cusp structure is defined by a fibration
\(\partial X\to Y\); when \(Y\) is a point, its vector fields are spanned
by \(x^2\partial_x\) and \(\partial_y\), and it is the cusp structure of
these notes. Their algebras \({}^\Phi\Is\), \({}^\Phi\As\),
\({}^\Phi\fA_{\partial,\sigma}\) are then \(\Is\), \(\As\), \(\Ads\), and
\cite[Thm.~5.3]{LauterMoroianu} with \(Y\) a point is
Propositions~\ref{prop:Is}, \ref{prop:As} and~\ref{prop:Ads}, convergence
included. The convergence rests on a criterion for spectral sequences with
three filtrations \cite[Thm.~B.4]{LauterMoroianu}, whose hypotheses are
verified in \cite[\S5]{LauterMoroianu}. We give the argument in the
setting of these notes, with three changes.
\begin{enumerate}
\item We work in the stages of the mixed complexes of
Lemma~\ref{ha:lem:mixed}, whose part (b) is now proved. In
\cite{LauterMoroianu} the chain spaces are those of Benameur and Nistor
\cite{BNcomplete}, which are
the chain spaces of Section~\ref{subsec:chains}, and a third filtration
index bounds the orders and poles of the individual tensor factors; a
stage plays that role here (Remark~\ref{su:rem:LM}).
\item The hypotheses that \cite{LauterMoroianu} call H4 and H5 (HKR with
bounds) are replaced by Lemma~\ref{ha:lem:stageE1}.
\item The degeneration at \(E^2\) of the order spectral sequence of the
suspended symbol algebra, hypothesis H3 of \cite{LauterMoroianu}, is
proved in full (Proposition~\ref{su:prop:degS}).
\cite[Prop.~5.1]{LauterMoroianu} only indicates the ingredients (central
multiplication operators, contractions and Lie derivatives along
derivations, and an exact sequence), and says they are enough; we use the
first two, organized through the weight of a single Lie derivative, and
do not need the exact sequence.
\end{enumerate}

\emph{The idea.} Lemma~\ref{ha:lem:WL} loses poles because it measures
them in the ordinary forms. Measured by the total power of \(x\) in a
chain, the primitive of an off-critical class loses at most one pole
(Lemma~\ref{su:lem:H1}). The remainder then has a part with one extra
pole. That part lives in the algebra graded by the power of \(x\), which
is \(\Ssym\otimes\C[x,x^{-1}]\) with \(\Ssym\) the suspended symbol algebra
and \(x\) central. If the order spectral sequence of that algebra
degenerates at \(E^2\), the extra part can be removed order by order with
no further loss of poles (Proposition~\ref{su:prop:step}). So all the
corrections stay in one stage.

\subsection{A convergence criterion}\label{su:subsec:criterion}

\subsubsection*{Setting}
Fix \(n\in\Z\). For \((L,M)\in\Z^2\), \(k\ge0\) and \(p\le M\),
\(\Lambda\le L\) let \(K^{p,\Lambda}_{L,M;k}\) be vector spaces, with
\(K^{p,\Lambda}_{L,M;k}\subset K^{p,\Lambda}_{L',M';k}\) for \(L\le L'\),
\(M\le M'\). The \emph{stages} are
\[
 S_{L,M;k}=\prod_{p\le M,\ \Lambda\le L}K^{p,\Lambda}_{L,M;k},
\]
and \(C_k=\bigcup_{L,M}S_{L,M;k}\). We write \(c^{p,\Lambda}\) for the
components of a chain \(c\), and call \(p\) the \emph{order} and
\(\Lambda\) the \emph{pole}. Suppose given linear maps
\(b_{m,l}:K^{p,\Lambda}_{L,M;k}\to K^{p-m-l,\Lambda-m}_{L,M;k-1}\) for
\(m,l\ge0\), compatible with the inclusions, and put
\[
 (bc)^{p,\Lambda}=\sum_{m,l\ge0}b_{m,l}\,c^{p+m+l,\Lambda+m}.
\]
For \(c\in S_{L,M}\) the sum is finite, since \(p+m+l\le M\) and
\(\Lambda+m\le L\). We assume \(b^2=0\). Then each stage is a subcomplex.
Let \(S^{p,\Lambda}_{L,M}\subset S_{L,M}\) consist of the chains whose
non-zero components have order \(\le p\) and pole \(\le\Lambda\); it is a
subcomplex. Let \(F_pC\) be the chains with components of order \(\le
p\), and \(c_p=(c^{p,\Lambda})_\Lambda\) the part of order \(p\). We write
\(b_0=b_{0,0}\) and \(b_1=b_{1,0}+b_{0,1}\). The spectral sequence of
\((C,F)\) has \(E^1_{p,k}=H_k\bigl(\prod_\Lambda K^{p,\Lambda}_{\bullet},
b_0\bigr)\) (in the direct limit) and \(d^1\) induced by \(b_1\). Let
\(G_\Lambda=\bigcup_{L,M}\prod_pK^{p,\Lambda}_{L,M}\) with the
differential \(b^\partial=\sum_lb_{0,l}\). It is a complex, since
\((b^\partial)^2\) is the part of \(b^2\) that preserves the pole. We
filter it by order.

\emph{Summability.} If \(c_j\in S_{L,M}\) for \(j\ge1\), in a fixed
stage, and every component \((p,\Lambda)\) is non-zero for only finitely
many \(c_j\), then \(\sum_jc_j\) is defined componentwise, and
\(b\sum_jc_j=\sum_jbc_j\), because each component of \(b\) involves
finitely many components. This holds, for example, if \(c_j\in
F_{p_j}\) with \(p_j\to-\infty\). It is the hypothesis H6 of
\cite{LauterMoroianu}.

\subsubsection*{Hypotheses}
\begin{enumerate}
\item[(H1)] \emph{(Off-critical primitives losing one pole.)} Let
\(p+1\le M-1\), \(\Lambda+1\le L-1\) and \(k\ne p+n\). If \(\alpha\in
S^{p,\Lambda}_{L,M;k}\) and \(b\alpha\in F_{p-2}\), there is
\(\beta\in S^{p+1,\Lambda+1}_{L,M;k+1}\) with \(\alpha-b\beta\in F_{p-1}\).
\item[(H3)] \emph{(Degeneration in the pole-graded complexes.)} For every
\(\Lambda\), the differentials \(d^r\), \(r\ge2\), of the spectral
sequence of the filtered complex \(G_\Lambda\) vanish.
\item[(H4)] \emph{(HKR in a stage.)} If \(p\le M-1\) and \(\Lambda\le
L-1\), the inclusion of \((K^{p,\Lambda}_{L,M;\bullet},b_0)\) into
\((\bigcup_{L',M'}K^{p,\Lambda}_{L',M';\bullet},b_0)\) is a
quasi-isomorphism.
\end{enumerate}
The numbering follows \cite[App.~B]{LauterMoroianu}; see
Remark~\ref{su:rem:LM}.

\begin{lemma}[Strictness]\label{su:lem:strict}
Let \((G,b)\) be a complex with an exhaustive increasing filtration
\(F_p\), whose spectral sequence has \(d^r=0\) for all \(r\ge2\). If
\(z\in F_p\) and \(z=bw\) with \(w\in F_{p+s}\) for some \(s\ge1\), then
\(z\in b(F_{p+1})+F_{p-1}\).
\end{lemma}

\begin{proof}
We use the standard description: with \(Z^r_p=\{x\in F_p:\ bx\in
F_{p-r}\}\), one has \(E^r_p=Z^r_p/(Z^{r-1}_{p-1}+bZ^{r-1}_{p+r-1})\) and
\(d^r[x]=[bx]\). Induct on \(s\); the case \(s=1\) is trivial. Let
\(s\ge2\). Then \(w\in Z^s_{p+s}\), and \(d^s[w]=0\) says \(bw\in
Z^{s-1}_{p-1}+bZ^{s-1}_{p+s-1}\). So \(z=z'+bw'\) with \(z'\in F_{p-1}\)
and \(w'\in F_{p+s-1}\). Since \(bw'=z-z'\in F_p\), the induction
hypothesis gives \(bw'\in b(F_{p+1})+F_{p-1}\).
\end{proof}

\begin{proposition}[Pole-preserving step; {\cite[Prop.~B.3]{LauterMoroianu}}]
\label{su:prop:step}
Assume (H1), (H3) and (H4). Let \(p+1\le M-1\), \(\Lambda+1\le L-1\),
\(k\ne p+n\), and let \(\alpha\in S^{p,\Lambda}_{L,M;k}\) with
\(b\alpha\in F_{p-2}\). Then there is \(\beta\in
S^{p+1,\Lambda+1}_{L,M;k+1}\) with
\[
 \alpha-b\beta\in S^{p-1,\Lambda}_{L,M;k}.
\]
The order goes down by one, and the pole bound \(\Lambda\) is restored.
\end{proposition}

\begin{proof}
For a chain \(c\) with poles \(\le\Lambda+1\) write \([c]_{\Lambda+1}=
(c^{p',\Lambda+1})_{p'}\in G_{\Lambda+1}\) for its part of pole
\(\Lambda+1\).

\emph{Step 1.} By (H1) there is \(\beta_1\in S^{p+1,\Lambda+1}_{L,M}\)
with \(R_1=\alpha-b\beta_1\in F_{p-1}\), so \(R_1\in
S^{p-1,\Lambda+1}_{L,M}\). We show by induction on \(j\ge1\) that there
are \(\beta_j\in S^{p+1,\Lambda+1}_{L,M}\) such that
\(R_j=\alpha-b\beta_j\) satisfies
\begin{equation}\label{su:eq:Rj}
 [R_j]_{\Lambda+1}\in F_{p-j},\qquad R_j-[R_j]_{\Lambda+1}\in
 S^{p-1,\Lambda}_{L,M},
\end{equation}
and \(\beta_{j+1}-\beta_j\in S^{p-j+1,\Lambda+1}_{L,M}\). For \(j=1\)
this holds.

\emph{Step 2.} Let \(z=[R_j]_{\Lambda+1}\). Since \(\alpha\) has no
component of pole \(\Lambda+1\), and \(b_{m,l}\) lowers the pole by \(m\),
\[
 z=-b^\partial[\beta_j]_{\Lambda+1}.
\]
So \(z\in G_{\Lambda+1}\cap F_{p-j}\) is the boundary of an element of
\(F_{p+1}\). By (H3) and Lemma~\ref{su:lem:strict}, applied in
\(G_{\Lambda+1}\), there are \(y\in G_{\Lambda+1}\cap F_{p-j+1}\) and
\(z'\in F_{p-j-1}\) with \(z=b^\partial y+z'\). The chain \(y\) may lie in
a larger stage. Comparing components of order \(p-j+1\) and \(p-j\),
\[
 b_0y_{p-j+1}=0,\qquad z_{p-j}=b_{0,1}y_{p-j+1}+b_0y_{p-j}.
\]
By (H4) at \((p-j+1,\Lambda+1)\) there is a \(b_0\)-cycle \(\xi\in
K^{p-j+1,\Lambda+1}_{L,M}\) with \(\xi-y_{p-j+1}=b_0v\) for some \(v\) in
a larger stage. Since \(b_0b_{0,1}+b_{0,1}b_0=0\) (a component of
\((b^\partial)^2=0\)),
\[
 z_{p-j}-b_{0,1}\xi=b_0\bigl(y_{p-j}+b_{0,1}v\bigr).
\]
The left side lies in \(K^{p-j,\Lambda+1}_{L,M}\), so by (H4) it equals
\(b_0\eta\) with \(\eta\in K^{p-j,\Lambda+1}_{L,M}\). Put
\(\beta_{j+1}=\beta_j+\xi+\eta\). The part of pole \(\Lambda+1\) of
\(R_{j+1}=R_j-b(\xi+\eta)\) is \(z-b^\partial(\xi+\eta)\). Its component
of order \(p-j+1\) is \(-b_0\xi=0\), and that of order \(p-j\) is
\(z_{p-j}-b_{0,1}\xi-b_0\eta=0\). The rest of \(b(\xi+\eta)\) comes from
the \(b_{m,l}\) with \(m\ge1\). It has pole \(\le\Lambda\) and order
\(\le p-j+1-m\le p-1\). This proves \eqref{su:eq:Rj} for \(j+1\).

\emph{Step 3.} The series \(\beta=\beta_1+\sum_{j\ge1}(\beta_{j+1}-
\beta_j)\) is summable in \(S^{p+1,\Lambda+1}_{L,M}\). The components of
\(R_j\) stabilize as \(j\to\infty\), so \(\alpha-b\beta=\lim R_j\)
componentwise. By \eqref{su:eq:Rj}, a component of pole \(\Lambda+1\)
vanishes in the limit, and the other components have order \(\le p-1\).
\end{proof}

\begin{theorem}[Convergence criterion; {\cite[Thm.~B.4]{LauterMoroianu}}]
\label{su:thm:criterion}
In the setting above, assume (H1), (H3) and (H4). Then:
\begin{enumerate}
\item \(E^2_{p,k}=0\) for \(p\ne k-n\), so \(d^r=0\) for \(r\ge2\);
\item every class in \(H_k(C)\) has a representative in \(F_{k-n}C_k\),
and \(\Theta[c]=[c_{k-n}]\in E^2_{k-n,k}\) is well defined;
\item \(\Theta:H_k(C)\to E^2_{k-n,k}\) is an isomorphism.
\end{enumerate}
\end{theorem}

\begin{proof}
Finitely many chains lie in a common stage, and a stage can be enlarged at
will. We therefore always choose stages with \(L\) and \(M\) at least two
larger than the poles and orders that occur, so that the margin
conditions of (H1) and Proposition~\ref{su:prop:step} hold.

(1) Let \(x\) be a \(b_0\)-cycle of order \(p\), degree \(k\), with
\(d^1[x]=0\), so \(b_1x=-b_0y\) with \(y\) of order \(p-1\). Then
\(c=x+y\) has \(bc\in F_{p-2}\). If \(p\ne k-n\), (H1) gives \(\beta\in
F_{p+1}\) with \(c-b\beta\in F_{p-1}\). The components of order \(p+1\)
and \(p\) give \(b_0\beta_{p+1}=0\) and \(x=b_1\beta_{p+1}+b_0\beta_p\),
so \([x]=d^1[\beta_{p+1}]=0\) in \(E^2\). Since \(d^r\) maps
\(E^r_{p,k}\) to \(E^r_{p-r,k-1}\), both can be non-zero only if
\(r=1\).

(2) A cycle \(c\in F_P\) with \(P>k-n\) satisfies \(bc=0\in F_{P-2}\),
and (H1) gives a homologous cycle in \(F_{P-1}\). After finitely many
steps \(c\in F_{k-n}\). For such a cycle, \(b_0c_{k-n}=0\) and
\(b_1c_{k-n}=-b_0c_{k-n-1}\), so \([c_{k-n}]\in E^2\). If moreover
\(c=be\) with \(e\in F_s\) and \(s\ge k+2-n\), then \(be\in F_{s-2}\) and
\(k+1\ne s+n\), so (H1) replaces \(e\) by \(e-b\beta\in F_{s-1}\), which
has the same boundary. Once \(s=k+1-n\), \(c_{k-n}=b_1e_{k+1-n}+
b_0e_{k-n}\) with \(b_0e_{k+1-n}=0\), so \([c_{k-n}]=d^1[e_{k+1-n}]=0\).

(3) \emph{Injectivity.} Let \(c\in F_{k-n}\) be a cycle with
\([c_{k-n}]=0\) in \(E^2\). Then \(c_{k-n}=b_1y'+b_0y''\), where \(y'\) is
a \(b_0\)-cycle of order \(k-n+1\) and \(y''\) has order \(k-n\). So
\(c^{(0)}=c-b(y'+y'')\in F_{k-n-1}\) is a cycle. Choose a stage
\(S_{L,M}\) containing it, with poles \(\le\Lambda\), \(\Lambda+1\le L-1\)
and \(k-n+1\le M-1\). For \(j\ge0\) the cycle \(c^{(j)}\) lies in
\(S^{k-n-1-j,\Lambda}_{L,M}\), and \(k\ne(k-n-1-j)+n\). So
Proposition~\ref{su:prop:step} gives \(\beta^{(j)}\in
S^{k-n-j,\Lambda+1}_{L,M}\) with \(c^{(j+1)}=c^{(j)}-b\beta^{(j)}\in
S^{k-n-2-j,\Lambda}_{L,M}\), again a cycle. The series
\(\sum_j\beta^{(j)}\) is summable in \(S_{L,M}\), and
\(c^{(0)}=b\sum_j\beta^{(j)}\), since the components of \(c^{(j)}\)
vanish in the limit. Hence \(c\) is a boundary.

\emph{Surjectivity.} Let \(x\) be a \(b_0\)-cycle of order \(k-n\) with
\(d^1[x]=0\), and \(y\) of order \(k-n-1\) with \(b_1x=-b_0y\). Choose a
stage as before containing \(x\) and \(y\), and put \(c_1=x+y\). Then
\(bc_1\in S^{k-n-2,\Lambda}_{L,M;k-1}\), and \(k-1\ne(k-n-2)+n\).
Inductively, suppose \(c_j\) has leading term \(x\) and \(bc_j\in
S^{k-n-1-j,\Lambda}_{L,M}\). Proposition~\ref{su:prop:step}, applied to
\(\alpha=bc_j\) (a cycle, of degree \(k-1\), off the critical line), gives
\(\beta_j\in S^{k-n-j,\Lambda+1}_{L,M}\) with \(b(c_j-\beta_j)\in
S^{k-n-2-j,\Lambda}_{L,M}\). Put \(c_{j+1}=c_j-\beta_j\). The limit
\(c=c_1-\sum_j\beta_j\) exists in \(S_{L,M}\), is a cycle, and has
leading term \(x\), because every \(\beta_j\) has order \(\le k-n-1\). So
\(\Theta[c]=[x]\).
\end{proof}

\begin{remark}[Comparison with {\cite[App.~B]{LauterMoroianu}}]
\label{su:rem:LM}
Lauter and Moroianu filter the chains of \cite{BNcomplete} by three
indices:
the total order \(i\), the total pole \(j\) (the sum of the pole orders of
the factors, which can be negative) and an index \(l\) bounding every
partial sum of the orders and of the pole orders of the factors. Their
chain space is \(\bigcup_N\fA_{N}^{\hot(k+1)}\), which is the chain space
of Section~\ref{subsec:chains}, so their Hochschild homology is the one of
these notes, and excision (Theorem~\ref{ha:thm:excision}) applies to it.
The index \(l\) makes their filtered pieces complete, which is what a
stage does here. Their hypotheses correspond to ours as follows. Our
(H1), which assumes only \(b\alpha\in F_{p-2}\), combines their H1, H2
(vanishing of \(E^2\) off the critical line) and part of H4, with the
same loss of one pole. Their H3, fed by their Corollary~B.8 (which uses
their Lemma~B.7), is our (H3) together with Lemma~\ref{su:lem:strict}.
Their H4 and H5 (HKR with bounds on \(l\), their Lemma~B.6) are replaced
by (H4), which Lemma~\ref{ha:lem:stageE1} gives in the stages; their
Lemma~B.7 also enters their H1 through their Proposition~5.2. Their H6
is summability, automatic in a stage. Proposition~\ref{su:prop:step} and
Theorem~\ref{su:thm:criterion} follow their Proposition~B.3 and
Theorem~B.4. The price of using stages is Lemma~\ref{ha:lem:mixed}(b);
the price of their index \(l\) is HKR with bounds on \(l\), for which
\cite{LauterMoroianu} give a short indication (an explicit homotopy to the
small complex), which we have not checked in detail.
\end{remark}

\subsection{The suspended symbol algebra}\label{su:subsec:S}

Let \(H=\partial X\), of dimension \(d=n-1\), and
\(W=T^*H\times\R_\tau\). Then \(Z=\Zb\) is the sphere bundle of \(W\), and
homogeneous functions of degree \(p\) on \(W\setminus0\) are \(r^p\)
times functions on \(Z\). Let
\[
 \Ssym=\Psis^{\Z}(H)/\Psis^{-\infty}(H),
\]
filtered by order. (\(\Ssym\) is not the Schwartz space.) We use the
following facts.
\begin{enumerate}
\item[(S1)] A quantization, given by (K1) of Background
fact~\ref{cc:bg:uniform} and Borel's lemma, identifies \(\Ssym_{\le m}\)
with \(\prod_{p\le m}\mathcal G^p(W)\), the products of homogeneous
functions on
\(W\setminus0\). The product is a star product with \(P_0(a,b)=ab\) and
\(P_1(a,b)-P_1(b,a)=-i\{a,b\}_H\), where \(\{\cdot,\cdot\}_H\) is the
Poisson bracket of \(T^*H\), with \(\tau\) as a parameter. Indeed, by
Proposition~\ref{cc:prop:indicial} the normal operator is multiplicative
and \(\sigma(N(A))=\sigma(A)|_{x=0}\), and at \(x=0\) the cusp bracket
\eqref{eq:poisson} reduces to \(\{\cdot,\cdot\}_H\).
\item[(S2)] Let \(\hat\tau\in\Ssym\) have indicial family \(\tau\) (it is
\(-D_t\)). It is central, because \(\In\) is multiplicative. Let
\(D=\mathrm{ad}_t\). It is a derivation of order \(-1\), with
\(\In(Da)=-i\partial_\tau\In(a)\) (Lemma~\ref{cc:lem:sus}), so
\(D\hat\tau=-i\). Put \(T=i\hat\tau D\). It is a derivation, since
\(\hat\tau\) is central, and \(\In(Ta)=\tau\partial_\tau\In(a)\). On
principal symbols \(D\) acts by \(-i\partial_\tau\) and \(T\) by the
vector field \(\tau\partial_\tau\).
\item[(S3)] With \(\Ssym\) as a filtered nuclear algebra
(\(\Ssym_N=\Ssym_{\le N}\)) and the chains of
Section~\ref{subsec:chains}, the spectral sequence of the order
filtration has \(E^1_{p,k}\) equal, via \(\chi\), to the homogeneous
\(k\)-forms of degree \(p\) on \(W\setminus0\)
(Corollary~\ref{ha:cor:homog} with \(R=\CI(Z)\)), and \(d^1=-i\delta_H\),
where \(\delta_H\) is the Koszul--Brylinski differential of the Poisson
structure of \(T^*H\) (Proposition~\ref{prop:d1} and (S1)). The function
\(\tau\) is a Casimir.
\end{enumerate}

Let \(\Omh\) be the forms \(\omega\) on \(W\setminus0\) with
\(\iota_{\partial_\tau}\omega=0\), i.e.\ those not containing \(d\tau\),
and \(\Omh^k{}_{(p)}\) those of degree \(k\) and homogeneity \(p\). Every
form is uniquely \(\omega=\alpha+d\tau\wedge\beta\) with \(\alpha,\beta\in
\Omh\).

\begin{lemma}[The \(E^2\) term of \(\Ssym\)]\label{su:lem:E2S}
\textup{(i)} \(\delta_H\) preserves \(\Omh\) and anticommutes with
\(d\tau\wedge\). Hence
\[
 E^2_{p,k}(\Ssym)=\mathfrak a_{p,k}\oplus d\tau\wedge\mathfrak
 a_{p-1,k-1},\qquad\mathfrak a_{p,k}=H_k\bigl(\Omh{}_{(\bullet)},
 \delta_H\bigr)_{p},
\]
the homology at \(\Omh^k{}_{(p)}\) of \(\delta_H:\Omh^k{}_{(p)}\to
\Omh^{k-1}{}_{(p-1)}\).

\textup{(ii)} The Lie derivative \(\mathcal L_{\tau\partial_\tau}\) acts
on \(\mathfrak a_{p,k}\) as multiplication by \(c=p+d-k\).
\end{lemma}

\begin{proof}
(i) The bivector \(\Pi_H\) of \(T^*H\) has no \(\partial_\tau\)
component and does not depend on \(\tau\). So \(\iota_{\Pi_H}\)
commutes with \(\iota_{\partial_\tau}\) and with \(d\tau\wedge\), and
\(\mathcal L_{\partial_\tau}\Pi_H=0\). From
\(\delta_H=\iota_{\Pi_H}d-d\,\iota_{\Pi_H}\) and Cartan's formula
\(\iota_{\partial_\tau}d+d\iota_{\partial_\tau}=\mathcal
L_{\partial_\tau}\) one finds \(\iota_{\partial_\tau}\delta_H=
-\delta_H\iota_{\partial_\tau}\) and \(\delta_H(d\tau\wedge\omega)=
-d\tau\wedge\delta_H\omega\). The first identity shows that \(\Omh\) is a
subcomplex. So \((E^1,d^1)\) is the direct sum of \((\Omh,-i\delta_H)\) and
\(d\tau\wedge(\Omh,-i\delta_H)\); the form \(d\tau\) has homogeneity
one.

(ii) On \(\Omh\) write \(d=d_F+d\tau\wedge\partial_\tau\), where \(d_F\)
is the exterior derivative of \(T^*H\), with \(\tau\) as a parameter, and
maps \(\Omh\) to itself. On \(\Omh\), \([\iota_{\Pi_H},d\tau\wedge
\partial_\tau]=0\), so \(\delta_H\) is the Koszul--Brylinski differential
of the symplectic manifold \(T^*H\), applied with \(\tau\) as a
parameter. Let \(*\) be the symplectic star of \(T^*H\), applied
pointwise to \(\Omh\). By Proposition~\ref{prop:star}, \(\delta_H=
(-1)^{k-1}*^{-1}d_F*\) on \(\Omh^k\), and \(*\) maps
\(\Omh^k{}_{(p)}\) onto \(\Omh^{2d-k}{}_{(p+d-k)}\), because the
symplectic form \(d\eta\wedge dy\) has homogeneity one. The vector field
\(\tau\partial_\tau\) preserves \(d\eta\wedge dy\) and \(\Omh\), so
\(\mathcal L_{\tau\partial_\tau}\) commutes with \(*\), \(d_F\) and
\(\delta_H\). Hence it suffices to show that \(\mathcal
L_{\tau\partial_\tau}=c\) on the \(d_F\)-cohomology of \(\Omh{}_{(c)}\).
Let \(R'=\eta\cdot\partial_\eta\). It has no \(\partial_\tau\) component
and does not depend on \(\tau\), so \(\iota_{R'}\) preserves \(\Omh\)
and \(\mathcal L_{R'}=d_F\iota_{R'}+\iota_{R'}d_F\) on \(\Omh\). The radial
field is \(R=R'+\tau\partial_\tau\), and \(\mathcal L_R=c\) on forms of
homogeneity \(c\). So \(\mathcal L_{\tau\partial_\tau}=c-(d_F\iota_{R'}+
\iota_{R'}d_F)\) on \(\Omh{}_{(c)}\).
\end{proof}

Lemma~\ref{su:lem:E2S} by itself does not decide where \(E^2(\Ssym)\)
lives, since \(d\tau\wedge\mathfrak a\) lies on the same line as
\(\mathfrak a\). But the classes \(\tau^c\,\omega_H^d/d!\), \(c\ge0\), where \(\omega_H\) is
the symplectic form of \(T^*H\) and \(d=\dim H\), survive and lie on the lines \(p-k=c-d\), so
\(E^2(\Ssym)\) is not concentrated on a line, and
Proposition~\ref{ha:prop:algconv} does not apply to \(\Ssym\)
(Remark~\ref{su:rem:HS}). What we need is that there are no higher
differentials.

\begin{proposition}[Degeneration for \(\Ssym\); compare
{\cite[Prop.~5.1]{LauterMoroianu}}]\label{su:prop:degS}
The differentials \(d^r\), \(r\ge2\), of the order spectral sequence of
\(\Ssym\) vanish.
\end{proposition}

\begin{proof}
\emph{Three operators on chains.} Let \(s_i\) insert \(\hat\tau\) after
the \(i\)-th slot, and
\begin{gather*}
 \varepsilon(a_0\otimes\dots\otimes a_k)=\sum_{i=0}^k(-1)^is_i(a_0\otimes
 \dots\otimes a_k),\qquad\iota=i_D,\\
 \mathcal L_T(a_0\otimes\dots\otimes a_k)=\sum_ja_0\otimes\dots\otimes
 Ta_j\otimes\dots\otimes a_k .
\end{gather*}
These are continuous on stages. They satisfy
\[
 b\varepsilon+\varepsilon b=0,\qquad b\iota+\iota b=0,\qquad
 \iota\varepsilon+\varepsilon\iota=D(\hat\tau)=-i,\qquad
 b\mathcal L_T=\mathcal L_Tb .
\]
The second is Lemma~\ref{lem:iD} and the fourth holds because \(T\) is a
derivation. For the first, the faces of \(bs_i\) that do not touch
\(\hat\tau\) cancel against \(\varepsilon b\). Among the faces that do,
\(\hat\tau\) multiplies the entry on its left or on its right. Merging
into \(a_{i+1}\) from \(s_i\) and from \(s_{i+1}\) gives the same chain,
because \(\hat\tau\) is central, with opposite signs. The two remaining
terms, \(a_0\hat\tau\otimes\cdots\) from \(s_0\) and
\(\hat\tau a_0\otimes\cdots\) from the last face of \(s_k\), also cancel.
For the third, \(i_Ds_0=D(\hat\tau)\cdot\) (the entry \(a_0D(\hat\tau)\)
in slot \(0\)), and \(i_Ds_i=s_{i-1}i_D\) for \(i\ge1\), so
\(\iota\varepsilon=-i-\varepsilon\iota\). (We also checked these
identities symbolically, with a free algebra and a central variable.)

The operators are filtered of degrees \(+1\), \(-1\), \(0\), since
\(\hat\tau\), \(D\), \(T\) have orders \(1\), \(-1\), \(0\). So they act on
every page \(E^r\), they commute or anticommute with \(d^r\), and the
identity \(\iota\varepsilon+\varepsilon\iota=-i\) holds on every page.
On \(E^1\), through \(\chi\): \(\varepsilon\) is \(d\tau\wedge\), since
\(\chi(s_i(a_0\otimes\dots\otimes a_k))=\frac{(-1)^i}{(k+1)!}\,d\tau\wedge
a_0\,da_1\wedge\dots\wedge da_k\) for the leading symbols; \(\iota\) is
\(-i\,\iota_{\partial_\tau}\), as one sees on the antisymmetrizations
\(\varphi(a_0;a_1,\dots,a_k)\), which represent all classes
(Lemma~\ref{ha:lem:chi}); and \(\mathcal L_T\) is \(\mathcal
L_{\tau\partial_\tau}\). On \(E^2\), in the decomposition of
Lemma~\ref{su:lem:E2S}, \(\iota\) vanishes on \(\mathfrak a\),
\(\varepsilon\) maps \(\mathfrak a_{p-1,k-1}\) onto the second summand of
\(E^2_{p,k}\), and \(\mathcal L_T=p+d-k\) on \(\mathfrak a_{p,k}\).

\emph{Induction.} Suppose \(d^2=\dots=d^{r-1}=0\) for some \(r\ge2\), so
that \(E^r=E^2\) with the same operators. Let \(a\in\mathfrak a_{p,k}\),
and write \(d^ra=a'+\varepsilon a''\) with \(a'\in\mathfrak
a_{p-r,k-1}\), \(a''\in\mathfrak a_{p-r-1,k-2}\). Since \(\iota a=0\), also
\(\iota d^ra=\mp d^r\iota a=0\). On the other hand \(\iota a'=0\) and
\(\iota\varepsilon a''=-ia''-\varepsilon\iota a''=-ia''\). So \(a''=0\).
Next, \(\mathcal L_T\) commutes with \(d^r\), and it acts by \(p+d-k\) on
\(a\) and by \(p-r+d-k+1\) on \(a'\). These differ, so \(a'=0\). Hence
\(d^r=0\) on \(\mathfrak a\), and \(d^r\varepsilon a=\mp\varepsilon
d^ra=0\). So \(d^r=0\).
\end{proof}

The weight \(\mathcal L_T\) separates the classes of \(\mathfrak a\) on
different lines \(p-k=\mathrm{const}\); the pair \(\varepsilon,\iota\)
separates \(\mathfrak a\) from \(d\tau\wedge\mathfrak a\), which carry
the same weight on adjacent lines. These are the ``exterior product by a
central function'' and ``contraction with a derivation'' of
\cite[Prop.~5.1]{LauterMoroianu}, with \(f=\tau\) and \(V=\partial_\tau\).

\begin{remark}[The size of \(E^2(\Ssym)\); sketch, not used]
\label{su:rem:HS}
By Lemma~\ref{su:lem:E2S}, \(\mathfrak a_{p,k}\) is the cohomology in
degree \(2d-k\) of the complex \((\Omh{}_{(c)},d_F)\), \(c=p+d-k\), the
leafwise de Rham complex of the foliation \(\{\tau=\mathrm{const}\}\) in
homogeneity \(c\). Passing to \(Z\), it becomes a complex of b-forms for
the hypersurface \(\{\tau=0\}=S^*H\subset Z\), weighted by \(s^{-c}\)
with \(s=\tau/|(\eta,\tau)|\). The Taylor coefficients at \(\{s=0\}\)
contribute only when \(c\in\{0,1,2,\dots\}\), and a Mayer--Vietoris
argument gives \(H^j(Z)\oplus H^{j-1}(S^*H)\) for \(c\ge0\) and
\(H^{j-d}(H)\oplus H^{j-d}(H)\) (one Thom class for each hemisphere
\(\pm\tau>0\)) for \(c<0\). In particular \(E^2(\Ssym)\) is non-zero on
every line \(p-k=\mathrm{const}\). We have not written out the details,
and nothing above uses them.
\end{remark}

\begin{remark}[The inner Euler derivation]\label{su:rem:euler}
In \(\Ads\) the element \(u=x^{-1}\hat\tau\) (pole one, order one) gives
an inner derivation with
\[
 [u,x^ja]=-i\,x^j\bigl(j+\tau\partial_\tau\bigr)a,\qquad a\in\Ssym,
\]
exactly, with \(\tau\partial_\tau\) acting on indicial families. Indeed,
by Proposition~\ref{cc:prop:taylor}, \([x^{-1},a]=\mathrm{ad}_t(a)\) and
\(x^{-1}\hat\tau x=\hat\tau-ix\), so \([u,x]=-ix\) and
\([u,a]=\mathrm{ad}_t(a)\hat\tau=-i\,Ta\). So \(\mathcal L_T\) of the proof
above is the restriction to \(\Ssym\) of an inner derivation of \(\Ads\),
up to \(x\partial_x\). We use \(T\) only through its action on
\(E^2(\Ssym)\), where it is semisimple. On chains it is not: on
homogeneous functions \(\tau\partial_\tau-\lambda\) is a first order
operator along the meridians of \(Z\) from \(\{\tau=0\}\) to
\(\{\eta=0\}\), with regular singular points at both ends. For generic
\(\lambda\) it is injective, but its range has infinite codimension:
solvability requires the vanishing of a Mellin transform along each
meridian. So the inner derivation alone does not give a chain-level
splitting by weight, and we do not use it in the proof.
\end{remark}

\subsection{Convergence for \texorpdfstring{\(\Ads\)}{Ads} and
\texorpdfstring{\(\As\)}{As}}\label{su:subsec:B10}

Let \(\fA=\Ads\), with the mixed complex \(C(\mathfrak B,\fA)\) of
Lemma~\ref{ha:lem:mixed}. By Proposition~\ref{cc:prop:taylor},
\(\fA=\Ssym\Lser{x}\), with coefficients of bounded order, and with the
product
\begin{equation}\label{su:eq:twist}
 (x^jg)(x^lh)=\sum_{i\ge0}\binom{l+i-1}{i}x^{j+l+i}\,\mathrm{ad}_t^i(g)\,h .
\end{equation}
By (S1), \(\fA_{L,M}=\prod_{J\ge-L,\,p\le M}x^J\mathcal G^p(W)\), and
\(\mathfrak B=\fA_{0,0}\). By (F3) a stage \(S_{L,M;k}=\fA_{L,M}\hot
\mathfrak B^{\hot k}\) is the product of the spaces
\(\mathcal G^{p_0}\hot\dots\hot\mathcal G^{p_k}\cong\CI(Z^{k+1})\) over
the multi-degrees \(((J_0,p_0),\dots,(J_k,p_k))\) with \(J_0\ge-L\),
\(p_0\le M\) and \(J_r\ge0\), \(p_r\le0\) for \(r\ge1\). Group them by
the total order \(\sum p_r\) and the total pole \(-\sum J_r\). For given
totals only finitely many multi-degrees occur. Let
\(K^{p,\Lambda}_{L,M;k}\) be the product of those of total order \(p\) and
total pole \(\Lambda\). Then \(S_{L,M;k}=\prod_{p\le M,\Lambda\le L}
K^{p,\Lambda}_{L,M;k}\). In \eqref{su:eq:twist} the term \(i\) raises the
power of \(x\) by \(i\) and lowers the order by at least \(i\). So the
Hochschild differential, which is continuous on stages, decomposes as
\(b=\sum_{m,l\ge0}b_{m,l}\), with \(b_{m,l}\) lowering the pole by \(m\)
and the order by \(m+l\). This is the setting of
Section~\ref{su:subsec:criterion}. Here \(b_0=b_{0,0}\) is the Hochschild
differential of the commutative algebra of leading terms, and
\(b^\partial=\sum_lb_{0,l}\) is that of
\[
 \Gsym=\Ssym\otimes\C[x,x^{-1}],
\]
with \(x\) central, because the terms of \eqref{su:eq:twist} with \(i=0\)
are \(x^{j+l}gh\). So \(G_\Lambda\) is the part of total \(x\)-degree
\(-\Lambda\) of the mixed complex \(C(\mathfrak B_\Gsym,\Gsym)\),
\(\mathfrak B_\Gsym=\Ssym_{\le0}[x]\).

Give \(x^a\) the weight \(a\) and \(dx\) the weight \(1\), and let
\(\Omega^k_{p,\Lambda}\) be the homogeneous formal \(k\)-forms of degree
\(p\) and weight \(-\Lambda\) on \(\cT\setminus0\) near
\(x=0\), that is, \(x^{-\Lambda}\alpha+x^{-\Lambda-1}dx\wedge\alpha'\) with
\(\alpha,\alpha'\) homogeneous forms on \(W\setminus0\).

\begin{lemma}[(H4) for \(\Ads\)]\label{su:lem:H4}
\(\chi\) identifies the \(b_0\)-homology of
\(\bigcup_{L,M}K^{p,\Lambda}_{L,M}\) with \(\Omega_{p,\Lambda}\). If \(p\le
M-1\) and \(\Lambda\le L-1\), the same holds for
\(K^{p,\Lambda}_{L,M}\). In particular (H4) holds.
\end{lemma}

\begin{proof}
The complex \(\bigl(\prod_\Lambda K^{p,\Lambda}_{L,M},b_0\bigr)\) is the
part of order \(p\) of the graded complex of the stage. By
Lemma~\ref{ha:lem:stageE1}, for \(p\le M-1\) its homology is, via
\(\chi\), the space of homogeneous formal forms of degree \(p\) of the
form \(\sum_{a\ge-L}x^a(\alpha_a+dx\wedge\alpha'_a)\). Since \(b_0\) and
\(\chi\) preserve the weight, this identification holds weight by weight.
The forms of weight \(-\Lambda\) in that space are
\(x^{-\Lambda}\alpha+x^{-\Lambda-1}dx\wedge\alpha'\), subject to
\(-\Lambda\ge-L\) and \(-\Lambda-1\ge-L\); for \(\Lambda\le L-1\) this is
all of \(\Omega_{p,\Lambda}\). The inclusions of stages induce the
identity on forms.
\end{proof}

\begin{lemma}[(H1) for \(\Ads\); {\cite[Prop.~5.2]{LauterMoroianu}}]
\label{su:lem:H1}
(H1) holds for \(\Ads\).
\end{lemma}

\begin{proof}
Let \(\alpha\in S^{p,\Lambda}_{L,M;k}\) with \(b\alpha\in F_{p-2}\),
\(k\ne p+n\), \(p+1\le M-1\), \(\Lambda+1\le L-1\). Then
\(b_0\alpha_p=0\) and \(b_1\alpha_p=-b_0\alpha_{p-1}\), so the form
\(\theta=\chi(\alpha_p)\) satisfies \(\delta\theta=0\), by
Proposition~\ref{prop:d1} (applied to mixed chains as in the proof of
Theorem~\ref{ha:thm:conv}). It is homogeneous of degree \(p\) and has
weights \(\ge-\Lambda\).

\emph{Weights and cusp forms.} In the coframe \(dx/x^2,dy,d\tau,d\eta\)
the cusp symplectic form \(\omega=d\tau\wedge dx/x^2+d\eta\wedge dy\) has
constant coefficients, so the symplectic star \(*\) maps a coframe
monomial to \(\pm\) a coframe monomial, and it preserves the space
\(x^{-\Lambda}\Omega_c\) of formal forms \(\sum_{a\ge-\Lambda}x^a
(\cdots)\) written in this coframe. So does \(\iota_R\). The monomial
\(dx/x^2\) has weight \(-1\), and the others weight \(0\). Hence a form of
weights \(\ge-\Lambda\) lies in \(x^{-\Lambda}\Omega_c\) (write
\(x^bdx=x^{b+2}dx/x^2\), with \(b+2\ge-\Lambda+1\)), and every form in
\(x^{-\Lambda}\Omega_c\) has weights \(\ge-\Lambda-1\).

\emph{The primitive.} As in the proof of Theorem~\ref{ha:thm:conv},
\(*\theta\) is closed and homogeneous of degree \(w=p+n-k\ne0\), and
\(\zeta=*^{-1}(\iota_R*\theta)/w\) satisfies \(\delta\zeta=(-1)^k\theta\)
(Proposition~\ref{prop:star}). It is homogeneous of degree \(p+1\), and it
lies in \(x^{-\Lambda}\Omega_c\), so it has weights \(\ge-\Lambda-1\).
By Lemma~\ref{su:lem:H4}, applied to each weight \(-\Lambda'\) with
\(\Lambda'\le\Lambda+1\le L-1\), there is a \(b_0\)-cycle \(\beta_0\in
\prod_{\Lambda'\le\Lambda+1}K^{p+1,\Lambda'}_{L,M}\) with
\(\chi(\beta_0)=(-1)^ki\zeta\), so that \(\chi(b_1\beta_0)=-i\delta
\chi(\beta_0)=\theta\). Then \(\alpha_p-b_1\beta_0\) is a \(b_0\)-cycle in
\(\prod_{\Lambda'\le\Lambda+1}K^{p,\Lambda'}_{L,M}\) with \(\chi=0\), so by
Lemma~\ref{su:lem:H4} it is \(b_0\gamma\) with \(\gamma\) in the same
product. Put \(\beta=\beta_0+\gamma\in S^{p+1,\Lambda+1}_{L,M}\). The
components of \(\alpha-b\beta\) of order \(p+1\) and \(p\) are
\(-b_0\beta_0=0\) and \(\alpha_p-b_1\beta_0-b_0\gamma=0\).
\end{proof}

This is where the one-pole loss enters: \(\zeta\) contains
\(dx/x^2\), of weight \(-1\). Measured in ordinary forms that is a loss
of two poles (Lemma~\ref{ha:lem:WL}). It is a loss of one in the total
power of \(x\), and Proposition~\ref{su:prop:step} recovers it.

\begin{lemma}[(H3) for \(\Ads\)]\label{su:lem:H3}
(H3) holds for \(\Ads\).
\end{lemma}

\begin{proof}
Let \(C(\Gsym)=\varinjlim_N\Gsym_N^{\hot(k+1)}\), with
\(\Gsym_N=\Ssym_{\le N}\otimes\mathrm{span}\{x^j:|j|\le N\}\), filtered by
order and graded by the total power of \(x\). The complex \(G_\Lambda\) is
a subcomplex of its part of degree \(-\Lambda\). Both have \(E^1\) equal,
via \(\chi\), to \(\Omega_{\bullet,\Lambda}\): for \(G_\Lambda\) by
Lemma~\ref{su:lem:H4}, and for \(C(\Gsym)\) by
Corollary~\ref{ha:cor:homog}, Proposition~\ref{ha:prop:kunneth} and
Theorem~\ref{ha:thm:poly}, the graded algebra being
\(\CI(Z)\otimes\C[r,r^{-1}]\otimes\C[x,x^{-1}]\). A filtered chain map
that is an isomorphism on \(E^1\) is an isomorphism on every \(E^r\). So
it suffices to prove that \(d^r=0\), \(r\ge2\), for \(C(\Gsym)\).

Let \(\Sigma\subset C(\C[x,x^{-1}])\) be the span of the cycles \(x^J\)
and \(x^{J-1}\otimes x\), \(J\in\Z\), a complex with zero differential.
The shuffle map gives a chain map \(C(\Ssym)\otimes\Sigma\to C(\Gsym)\)
(Background fact~\ref{ha:bg:EZ}). It preserves the order and the power of
\(x\). The spectral sequence of \(C(\Ssym)\otimes\Sigma\), filtered by
the order of the first factor, is \(E^r(\Ssym)\otimes\Sigma\). On
\(E^1\) the map is \(\theta\otimes x^J\mapsto x^J\theta\),
\(\theta\otimes x^{J-1}\otimes x\mapsto\theta\wedge x^{J-1}dx\)
(Proposition~\ref{ha:prop:kunneth}), which is an isomorphism onto the
\(E^1\) of \(C(\Gsym)\). So the two spectral sequences agree from \(E^1\)
on, and Proposition~\ref{su:prop:degS} gives \(d^r=0\) for \(r\ge2\).
\end{proof}

\begin{theorem}[(B10) for all \(n\)]\label{su:thm:B10}
For every \(n\ge1\), the leading-symbol maps
\[
 \Theta:\HH_k(\Ads)\to E^2_{k-n,k}(\Ads),\qquad
 \Theta:\HH_k(\As)\to E^2_{k-n,k}(\As)
\]
are isomorphisms, and the order spectral sequences degenerate at \(E^2\).
Consequently Propositions~\ref{prop:As}, \ref{prop:Ads}
and~\ref{prop:boundaryLES} hold without quoted input.
\end{theorem}

\begin{proof}
For \(\Ads\), Lemmas~\ref{su:lem:H4}, \ref{su:lem:H1} and~\ref{su:lem:H3}
verify the hypotheses of Theorem~\ref{su:thm:criterion} for the mixed
complex \(C(\mathfrak B,\Ads)\), with \(n=\dim X\). That theorem gives the
isomorphism for the homology of the mixed complex, and
Lemma~\ref{ha:lem:mixed}(b) identifies this homology with
\(\HH(\Ads)\). Its \(E^2\) term is that of Proposition~\ref{prop:Ads}, and
\(\Theta\) is the map of Proposition~\ref{ha:prop:laurent}(1). The
statement for \(\As\) follows from Proposition~\ref{ha:prop:equiv}, which
also gives Proposition~\ref{prop:boundaryLES} with the constant
\((-1)^ki\).
\end{proof}

For \(n=1\) the algebra \(\Ssym\) is commutative and
Proposition~\ref{su:prop:degS} is trivial; Theorem~\ref{su:thm:B10} then
gives a second proof of Proposition~\ref{ha:prop:xadic}. The obstruction
of Example~\ref{ha:ex:obstruction} does not arise: it concerns classes at
the lowest pole of a stage, and the margins \(\Lambda+1\le L-1\) keep all
chains of the proof away from it.
